\scanpage{150}
\[
\tag{57}
\begin{pmatrix}
 A+2\unsure{\nu}C&B+2\unsure{\nu}D\\
 C&D
\end{pmatrix}
=\varepsilon\begin{pmatrix}R&S\\T\,(=S)&U\end{pmatrix}.
\]

Mais cette condition implique d'abord que la matrice du premier membre est symétrique (puisque $S=T$) i.e.
\[
2\nu D=C-B
\]
condition qui équivaut aussi, par (54), à
\[
\tag{58}
\inserted{2\nu D=1-D\quad\text{i.e.}}\qquad
\boxed{\mu D=1}\quad
\deleted{\mu D=\unsure{1}\;\illegible{1}\qquad\unsure{\mu}=\unsure{1}}\quad
\text{i.e.}\quad D=1/\mu,\quad\mu=1/D
\]
\inserted{où $\mu=2\nu+1$,}
ce qui implique déjà \uline{$\mu$ et $D$ inversibles}.

\deleted{D'autre part, on doit avoir}
Symétriquement, si on écrit (57) sous la forme
\[
\varepsilon\begin{pmatrix}A&B\\C&D\end{pmatrix}
=\unsure{\varepsilon}\begin{pmatrix}R-2\nu T&S-2\nu U\\T&U\end{pmatrix},
\]
la condition $C+D-B=1$ devient, sur les éléments de la matrice $\varphi$
\[
T+U-(S-2\nu U)=\varepsilon\ ;
\]
i.e.
\[
(2\nu+1)U=\varepsilon\quad\text{i.e.}
\]
\[
\tag{59}
\mu U=\varepsilon\quad\text{i.e.}\quad U=\varepsilon/\mu,
\quad\inserted{\mu=\varepsilon/U}
\]
ce qui équivaut aussi, tenant compte de (58), à la relation
\[
\tag{60}
D=\varepsilon U\quad\text{ou}\quad U=\varepsilon D
\]
qui n'est autre que l'égalité des termes des lignes inférieures des matrices (57).
\margin{149}

\scanpage{151}
\mspage{365}
On peut expliciter les relations (57), comme \deleted{\illegible{1}} quatre relations sur les $9$ variables $a,b,c,d,r,s,t,u,\mu$ en plus des relations (48), (49) ;
\[
\tag{61}
\left\{\begin{gathered}
\boxed{\begin{aligned}
D&=\varepsilon U\\
\mu D&=1\\
C&=\varepsilon T
\end{aligned}}
\quad\begin{aligned}
\deleted{\mu\unsure{U}=\unsure{\varepsilon}}{i.e.}\quad\deleted{\unsure{D=1/\mu}}\quad\mu U=\varepsilon\\
&
\end{aligned}\\
A=\varepsilon\bigl(R-\deleted{\unsure{2\nu}}\underbrace{(\mu-1)}_{2\nu}\underbrace{T}_{C}\bigr)
\end{gathered}\right.
\]
\margin{\uline{NB} cette dernière est conséquence des autres, compte tenu de $C+D-B=1$ et $AD-BC=1$}
soit -- \deleted{ou} \inserted{encore}
\[
\tag{62}
\left\{\begin{gathered}
\boxed{\begin{aligned}
\mu(c^2+d^2+cd)&=\delta\\
\unsure{c}^2+d^2+cd&=\varepsilon\delta(t^2+u^2)\\
\uline{ac}+\uline{bd}+\uline{ad}-(c^2+d^2+cd)&=\delta\varepsilon(rt+su)
\end{aligned}}\\
a^2+b^2+ab-(ac+bd+bc)=\delta\varepsilon\bigl(r^2+s^2-(\mu-1)(rt+su)\bigr)
\end{gathered}\right.
\]
\margin{conséquence des précédentes}

Quand les signes $\varepsilon,\delta$ sont fixés, on trouve \inserted{avec (48), (49)} un système de $6$ équations en $9$ variables $a,b,c,d,r,s,t,u,\mu$ -- on doit pouvoir vérifier que le schéma (sur $\operatorname{Spec}\mathbb Z$) des solutions de ces équations est de dim. relative $3$, et de dim relative $1$ sur les ``facteurs'' $(a,b,c,d)$ resp $(r,s,t,u)$. On a donc un point adélique de ce schéma absolu -- un pt à coefficients dans $\widehat{\mathbb Z}$\ldots
\margin{non, il est de dim. rel totale $4$}

\uline{Théorème} Le multiplicateur $\mu$ est déterminé en fonction de $\unsure{\xi|\zeta}\inserted{=\begin{pmatrix}A&B\\C&D\end{pmatrix}}$ (considéré comme élément de $\operatorname{Sl}(2,\widehat{\mathbb Z})$), par la relation
\[
\mu D=1\qquad(58)
\]
\unsure{\textup{(+)}}
et il est déterminé \uline{au signe près} par $\varphi=\begin{pmatrix}R&S\\T&U\end{pmatrix}$ par la relation
\[
\mu U=\varepsilon\qquad(59)\qquad\text{où }\varepsilon\in\{\pm1\},
\]
\margin{150}

\scanpage{152}
et où $\varepsilon$ est le signe qui intervient dans les formules
\[
\left\{\begin{aligned}
\unsure{\xi|\zeta}&=\varepsilon l_0^\nu\varphi\qquad(56)\\
g\rho(g)\rho^2(g)&=\varepsilon\qquad(45).\\
\inserted{g}&\inserted{=\varepsilon\alpha\rho(\alpha)^{-1}}
\end{aligned}\right.
\]
On a $\varepsilon=1$ sss $\mu\equiv1\ (4)$, $\varepsilon=-1$ sss $\mu\equiv-1\ (4)$.
\inserted{i.e. $\nu\equiv0\ (2)$}
\inserted{i.e. $\nu\equiv1\ (2)$}

\margin{en d'autres termes, on a $\mu\equiv\varepsilon\ (4)$ car $U=t^2+u^2\equiv1\ (4)$ donc $1/U\equiv(1)\ (4)$ donc $\mu=\varepsilon/U\equiv\varepsilon\ (4)$. Ce serait le moment d'inclure \unsure{aussi la|ici la} relation $\mu\equiv\delta\ (3)$ cf. avec la démonstration idoine\ldots}

Pour le prouver, on se rappelle que
\[
U=t^2+u^2
\]
notons que les seules valeurs mod $4$ prises par $t^2$ et $u^2$ sont $0$ et $1$, donc celles prises par $t^2+u^2$ sont $0,1$ et $2$, et la seule unité parmi celles-ci est $1$. Donc \deleted{\unsure{$\mu$}} $1/\mu$ peut se mettre sous la forme $\varepsilon(\unsure{t^2}+u^2)$ mod $4$ sss $1/\mu\equiv\varepsilon\ (4)$ ou encore $\mu\equiv\varepsilon\ (4)$ cqfd.

\deleted{Ainsi} On avait vu de façon toute analogue, grâce à la formule $\mu D=1\quad(58)$ et
\[
D=\delta(c^2+d^2+cd)
\]
que $\mu\equiv\delta\ (3)$ i.e. $\delta=+1$ (resp $-1$) sss $\mu\equiv1\ (3)$ (resp. $\mu\equiv-1\ \mathrm{mod}\ 3$). On \deleted{\unsure{pourrait donc}} notera, pour \unsure{tt|les} nb premiers $p$, par $\unsure{\varepsilon_p|\beta_p}$ la valeur en $\mu$ du ``caractère quadratique'' relatif à $p$, égal à $1$ sss $\mu_p\in(\mathbb Z_p^*)^2$. On aura donc
\[
\tag{63}
\left\{\begin{aligned}
\delta&=ad-bc=\unsure{\varepsilon_3|\beta_3}(\mu)=\varepsilon_3\\
\varepsilon&=D/U=\unsure{\varepsilon_2|\beta_2}(\mu)=\varepsilon_2,
\end{aligned}\right.
\]
et la connaissance de $\varepsilon_3$ (resp. $\varepsilon_2$, resp. des deux) équivaut à celle de $\mu$ mod $\unsure{3|2}$ (resp. mod $4$, resp. mod $12$).
\margin{$(\mathbb Z/12\mathbb Z)^*\simeq(\mathbb Z/4\mathbb Z)^*\times(\mathbb Z/3\mathbb Z)^*\simeq\{\unsure{\pm1}\}\times\{\pm1\}$}
\margin{151}

\scanpage{153}
\mspage{366}
Nous allons aussi calculer \deleted{\illegible{1}} (cf formules (45))
\[
\alpha=(l_0^\nu f)^{-1}\alpha_0=f^{-1}l_0^{-\nu}\alpha_0
\]
et
\[
g=\varepsilon\alpha\rho(\alpha)^{-1}
=\underbrace{\sigma_\infty(f)^{-1}}_{{}^t f}\,l_1^\nu\rho(f)
=(l_0^\nu f)^{-1}\unsure{\xi|\zeta}\,\rho(l_0^\nu f).
\]
On trouve
\[
\tag{64}
\begin{aligned}
\alpha&=\underbrace{\begin{pmatrix}u&-s\\-t&r\end{pmatrix}}_{f^{-1}}
\underbrace{\begin{pmatrix}1&2\nu\\0&1\end{pmatrix}
\begin{pmatrix}a&b\\c&d\end{pmatrix}}_{\begin{pmatrix}a+2\nu c&b+2\nu d\\c&d\end{pmatrix}}\\
&=\begin{pmatrix}u(a+2\nu c)-cs&u(b+2\nu d)-ds\\-t(a+2\nu c)+cr&-t(b+2\nu d)+dr\end{pmatrix},
\end{aligned}
\]
\margin{\uline{NB} Contrairement à $\alpha_0$ et $f$, et à $\unsure{\xi|\zeta},\varphi$ qui s'en déduisent, $\alpha$ \deleted{\illegible{1}} partage avec $g$ l'avantage et l'intérêt d'être déterminé sans ambiguïté (au signe près) par un automorphisme $u$ de $\widehat\Pi_{0,3}$, alors que $f$ etc sont \illegible{1} \unsure{choix d'ambiguïtés} $f\mapsto l_0^n f$ etc\ldots}
\[
\begin{aligned}
g&=\underbrace{\begin{pmatrix}r&t\\s&u\end{pmatrix}}_{\sigma_\infty(f^{-1})={}^t f}
\underbrace{\begin{pmatrix}1&0\\2\nu&1\end{pmatrix}}_{\unsure{l_1^\nu|l_0^\nu}}
\underbrace{\begin{pmatrix}t+u&-t\\\unsure{t+u-r-s}&\unsure{r-t}\end{pmatrix}}_{\rho(f)\ (\text{cf \S43 formules }(77))}\\
&=\begin{pmatrix}r&t\\s&u\end{pmatrix}
\begin{pmatrix}t+u&-t\\\unsure{\mu(t+u)-r-s}&-\mu t+r\end{pmatrix}\\
&=\begin{pmatrix}1+\mu t(t+u)&-\mu t^2\\-1+\mu u(t+u)&1-\mu ut\end{pmatrix},
\end{aligned}
\]
soit
\[
\tag{65}
g=\begin{pmatrix}1+\mu t(t+u)&-\mu t^2\\-1+\mu u(t+u)&1-\mu ut\end{pmatrix}.
\]

Il apparaît que $g$ s'exprime simplement en termes de $f$ i.e. de $r,s,t,u$ directement, et non en termes de $\alpha_0$ ou $\unsure{\xi|\zeta}$ i.e. via les $a,b,c,d$ ou $A,B,C,D$. Par contre, pour exprimer $\alpha$ (64), il faut faire intervenir à la fois \inserted{les} $r,s,t,u$ et les $a,b,c,d$.

Il était naturel, en relation avec l'hom canonique
\[
\mathcal N_{1,1}\longrightarrow\operatorname{Gl}(2,\widehat{\mathbb Z})
\]
(cf \S39, (11)) de chercher une matrice \deleted{\illegible{1}}, que par abus de langage je vais désigner par la même lettre, $u$, que l'automorphisme envisagé de $\widehat\Pi_{0,3}$ dont on était parti (tout comme je désigne par la même lettre les éléments $g,f$ etc de $\widehat\Pi_{0,3}$ ou de $\unsure{\widehat{\mathfrak S}_{0,3}|\widehat{\mathcal S}_{0,3}}$, et leurs images dans $\operatorname{Gl}(2,\widehat{\mathbb Z})'$, voire même des éléments correspondants dans $\operatorname{Gl}(2,\widehat{\mathbb Z})$ lui-$\overline{\mathrm m}$),
\margin{152}

\scanpage{154}
\[
u\in\operatorname{Gl}(2,\widehat{\mathbb Z})
\]
\deleted{\unsure{satisf}} de déterminant égal à $\mu$, et induisant sur $\operatorname{Sl}(2,\widehat{\mathbb Z})'$, par autom. intérieure, le même autom. que \inserted{le} $u$ \inserted{initial} (par passage au quotient), ce qui s'explicite par les formules
\[
\tag{66}\det u=\mu
\]
\[
\tag{67}u\sigma_\infty u^{-1}=\varepsilon'\sigma_\infty\qquad\varepsilon'\in\{\pm1\}
\]
\[
\tag{68}u\rho u^{-1}=\varepsilon''g\rho\qquad\varepsilon''\in\{\pm1\}
\]
En fait, \deleted{dans} (68) \deleted{on \unsure{devrait avoir}} \inserted{s'écrit}
\[
[u,\rho]=u\rho(u)^{-1}=\varepsilon''g
\]
et appliquant aux deux membres l'opération $x\mapsto x\rho(x)\rho^2(x)$, et tenant compte que $g\rho(g)\rho^2(g)=\varepsilon_2$, \deleted{on trouve} \inserted{et $x\rho(x)\rho^2(x)=1$ pour $x=u\rho(u)^{-1}$,} on trouve
\[
\tag{68 bis}\varepsilon''=\varepsilon_2
\]
et (68) équivaut donc à
\[
u\rho(u)^{-1}=\varepsilon_2g=\alpha\rho(\alpha)^{-1}
\]
\deleted{\unsure{d'où}} i.e. à
\[
\tag{69}u=\alpha u_0\quad\text{avec }[u_0,\rho]=1
\quad\bigl(\text{i.e. }u_0\rho(u_0)^{-1}=1\bigr).
\]

Nous avons déjà remarqué que (67), (68) déterminent $u$ seulement modulo multiplication par un scalaire $\lambda\in\widehat{\mathbb Z}^*$, et compte tenu de (66), $u$ est déterminé seulement modulo $\lambda$ satisfaisant $\lambda^2=1$ i.e. modulo ${}_2\widehat{\mathbb Z}^*\simeq\{\pm1\}^{\mathbb P}$ (où $\mathbb P$ est l'ens des nbs premiers), mais on espère néanmoins tomber sur un choix canonique.

Dans mes premiers calculs, j'avais pas encore \unsure{dégagé|dégoté} (69), où la condition $u_0$ commute à $\rho$ s'exprimait sous la forme
\[
\tag{70}u_0=p\,\unsure{\mathrm{id}}+q\rho=\begin{pmatrix}p&q\\-q&p+q\end{pmatrix}\quad(p,q\in\widehat{\mathbb Z}),
\]
\margin{153}

\scanpage{155}
\mspage{367}

de sorte que
\[
\det u_0=p^2+pq+q^2
\tag{71}
\]
et que (66) s'écrit (compte tenu de (69), \unsure{où}
$\det\unsure{\alpha}=\delta=\varepsilon_3$)
\[
p^2+pq+q^2=\mu\varepsilon_3
\tag{\unsure{72}}
\]
ce qui, compte tenu de la relation
\[
c^2+\unsure{cd}+d^2=\varepsilon_3/\mu
\qquad(52\text{ et }58)
\]
suggérait des relations de la forme
\[
p=\deleted{\varepsilon}\mu c,\ q=\mu d,\quad
\text{ou}\quad p=\mu d,\ q=\mu c.
\]
En fait, exprimant (67) comme équation en $p,q$
(s'ajoutant \unsure{à|:} 72), \unsure{je trouvais}
(au signe près bien \unsure{sûr}
qui ici est inessentiel)
\[
p=\mu d,\ q=\mu c\qquad\text{i.e.}\qquad u_0=\mu(c\rho+d)
\tag{73}
\]
mais en me plaçant dans le cas $\varepsilon'=+1$ exclusivement.
(J'avais cru conclure, un peu trop hâtivement, que le
signe $\varepsilon'$ était toujours $+1$, avant de m'apercevoir
que ce n'était guère compatible avec le corps des
classes, et de subodorer que l'on devrait avoir $\varepsilon'=+1$
sss $p\equiv1\ (4)$.) En faisant ce détour assez pesant
par les $(a,b,c,d)$, je suis quand même arrivé, après
bien des péripéties (erreurs de calcul etc) à la
formule où il n'y avait plus les $a,b,c,d$ du tout,
et d'une simplicité déconcertante :
\[
\uuline{u}=\mu(t\sigma_\infty+u)
=\mu\begin{pmatrix}u&t\\-t&u\end{pmatrix},
\tag{74}
\]
qui correspond manifestement au cas $\unsure{\varepsilon'_2|\varepsilon'}=1$, i.e.
celui où $\uline{u}$ \uline{commute} à $\sigma_\infty$. Comme

\margin{($\Longleftrightarrow$ sans doute, $\varepsilon_2=1$)}
\margin{Je souligne la lettre $u$ qui désigne une matrice, pour
la distinguer du $u$ qui désigne un scalaire.}
\margin{154}

\scanpage{156}

\[
\det(t\sigma_\infty+u)=t^2+u^2=\varepsilon_2/\mu
\quad(=1/\mu\text{ si }\varepsilon_2=+1),
\]
on aura bien, dans le cas $\varepsilon_2=1$, \deleted{\unsure{donc}}
\[
\det\uline{u}=\mu^2\det(t\sigma_\infty+u)=\mu^2\,1/\mu=\mu
\]
i.e., (66). D'autre part, on aura
\[
\begin{aligned}
[u,\rho]&=[t\sigma_\infty+u,\rho]\\
&=
\underbrace{
\begin{pmatrix}u&t\\-t&u\end{pmatrix}
\begin{pmatrix}0&1\\-1&1\end{pmatrix}
}_{\begin{pmatrix}-t&u+t\\-u&u-t\end{pmatrix}}
\inserted{\mu}
\underbrace{
\begin{pmatrix}u&-t\\t&u\end{pmatrix}
\begin{pmatrix}1&-1\\1&0\end{pmatrix}
}_{\begin{pmatrix}u-t&-u\\u+t&-t\end{pmatrix}}.
\\
&=\mu\begin{pmatrix}
(u+t)^2-ut+t^2&-t^2\\
ut-t^2&u^2+t^2-ut
\end{pmatrix}
=\mu\begin{pmatrix}
\dfrac1\mu+\unsure{t(u+t)}&-t^2\\
t(u-t)&\dfrac1\mu-ut
\end{pmatrix}
\\
&=\begin{pmatrix}
1+\mu t(u+t)&-\mu t^2\\
\mu t(u\deleted{\unsure{t}}-t^{\deleted{\unsure{2}}})&1-\mu ut
\end{pmatrix},
\end{aligned}
\]
\margin{On pourrait aussi \unsure{appliquer} les formules §43 (78) ---}
et en comparant avec (65) on trouve bien
\[
[u,\rho]=g
\]
compte tenu que
\[
\mu(u^2+t^2)=\varepsilon_2=1
\]
donc
\[
-1+\mu u(t+u)=\underbrace{-1+\mu u^2}_{-\mu t^2}+\mu ut
=\mu t(u-t)\ \unsure{-!|!}
\]
J'aimerais maintenant, sans me lancer dans de
nouveaux calculs, trouver "au pif" une expression
pour $u$, dans le \deleted{\unsure{cas}} où $\varepsilon_2=-1$, et où on s'attend
à trouver dans (67) le signe $\varepsilon'=-1$ --- i.e. que
$u$ soit une \uline{anti}\unsure{similitude}\illegible{1}
\[
\uline{u}\,{}^t\uline{u}=-\det\uline{u}
\]
Mais \deleted{\unsure{et}} la matrice
\[
\tau_\infty=\begin{pmatrix}-1&0\\0&1\end{pmatrix}
\]
\margin{(c'est \unsure{$=$} une réflexion \unsure{orthog.})}
est le \unsure{cas-type} d'une telle antisimilitude, et je
prévois que la matrice
\[
\uline{u}=\mu\tau_\infty(t\sigma_\infty+u)
=\mu\begin{pmatrix}\unsure{1|-1}&0\\0&-1\end{pmatrix}
\begin{pmatrix}u&t\\-t&u\end{pmatrix}
=\mu\begin{pmatrix}\unsure{u|-u}&\unsure{t|-t}\\-t&-u\end{pmatrix}
\]
\margin{155}

\scanpage{157}
\mspage{368}

conviendra. On aura, bien sûr
\[
\begin{aligned}
u\,{}^tu
&=\mu^2\begin{pmatrix}u&t\\t&-u\end{pmatrix}
\begin{pmatrix}u&t\\t&-u\end{pmatrix}\\
&=\mu^2\begin{pmatrix}u^2+t^2&0\\0&u^2+t^2\end{pmatrix}
=\mu^2\begin{pmatrix}-1/\mu&0\\0&-1/\mu\end{pmatrix}
\end{aligned}
\]
i.e.
\[
u\,{}^tu=-\mu.\mathrm{Id}
\]
et
\[
\det u=\mu^2(-u^2-t^2)=\mu^2(1/\mu)=\mu
\]
compte tenu que maintenant
\[
u^2+t^2=\varepsilon_2/\mu=-\frac1\mu
\qquad\text{puisque }\varepsilon_2=-1.
\]
Enfin pour vérifier (68) calculons
\[
\begin{aligned}
[u,\rho]
&=\left[\begin{pmatrix}\unsure{u}&t\\t&-u\end{pmatrix},\rho\right]\\
&=\underbrace{
\begin{pmatrix}u&t\\t&-u\end{pmatrix}
\begin{pmatrix}0&1\\-1&1\end{pmatrix}
}_{\begin{pmatrix}-t&u+t\\u&-u+t\end{pmatrix}}
\unsure{\dfrac1\mu|\mu}
\underbrace{
\begin{pmatrix}-u&-t\\-t&u\end{pmatrix}
\begin{pmatrix}1&-1\\1&0\end{pmatrix}
}_{\begin{pmatrix}-u-t&\unsure{u|-u}\\u-t&t\end{pmatrix}}
\\
&=\mu\begin{pmatrix}
u(u+t)&t^2\\
\unsure{-(u^2+t^2)}-u(u-t)&u^2+t^2-ut
\end{pmatrix}
=\begin{pmatrix}
\mu u(u+t)&\mu t^2\\
1-\mu u(u-t)&-1-\mu ut
\end{pmatrix}
\end{aligned}
\]
(compte tenu que $\mu(u^2+t^2)=-1$), \deleted{\unsure{d'où|donc}} a-t-on bien
\[
[u,\rho]=-g
\]
tenu
compte de la formule (65) qui explicite $g$ en
termes de $r,s,t,u$, on trouve que cette formule
\unsure{prévue|précise} équivaut à
\[
2\mu ut=0\qquad\text{i.e.}\qquad ut=0
\]
décidément, ça déconne ! Mais je vois que si on
remplace dans la définition de $u$ $t$ par $-t$,
i.e. si on définit
\[
\uline{u}=\mu(t\sigma_\infty+u)\tau_\infty
=\mu\begin{pmatrix}u&t\\-t&u\end{pmatrix}
\begin{pmatrix}1&0\\0&-1\end{pmatrix}
=\mu\begin{pmatrix}u&-t\\-t&-u\end{pmatrix}
\]
(au lieu de $\uline{u}=\mu\begin{pmatrix}u&t\\t&-u\end{pmatrix}$),
alors $\uline{u}$ est \unsure{tjs}
une antisimilitude de déterminant $\mu$, et de
plus maintenant
\[
[u,\rho]=\begin{pmatrix}
\mu u(u-t)&\mu t^2\\
1-\mu u(u+t)&-1+\mu ut
\end{pmatrix}=-g
\]
\margin{156}

\scanpage{158}

ce qui se réduit en effet à vérifier
\[
\mu u(u-t)=-(1+\mu t(t+u))
\]
i.e.
\[
\mu[u^2-\deleted{ut}+t^2+\deleted{tu}]=-1
\qquad\text{O.K.}
\]
En résumé, on trouve

\uline{Théorème} Soit $\alpha_2\in\mathbb Z/2\mathbb Z$ défini par
$(-1)^{\alpha_2}=\varepsilon_2$
(donc $\alpha_2\equiv0\ (2)$ sss $\mu\equiv1\ (4)$,
$\alpha_2\equiv1\ (2)$ sss
$\mu\equiv-1\ (4)$). Posons
\[
\uline{u}=\mu(t\sigma_\infty+u)\tau_\infty^{\alpha_2}
=\mu\begin{pmatrix}u&t\\-t&u\end{pmatrix}
\begin{pmatrix}1&0\\0&\varepsilon_2\end{pmatrix}^{\unsure{\alpha_2}},
\tag{75}
\]
\margin{$\displaystyle
=\begin{pmatrix}u&\varepsilon_2t\\-t&\varepsilon_2u\end{pmatrix}$}
\margin{\deleted{$\displaystyle
=\left\{\begin{array}{ll}
\mu\begin{pmatrix}u&t\\-t&u\end{pmatrix}&\alpha_2=0\\
\mu\begin{pmatrix}u&-t\\-t&-u\end{pmatrix}&\alpha_2=1
\end{array}\right.$}}
On a alors les formules
\[
\det\uline{u}=\mu
\tag{76}
\]
\[
[\uline{u},\sigma_\infty]=\unsure{-|\cdot}\varepsilon_2.\mathrm{id}
\qquad\text{i.e.}\quad\uline{u}(\sigma_\infty)=\varepsilon_2.\sigma_\infty
\tag{77}
\]
\[
[\uline{u},\rho]=\varepsilon_2.g
\qquad\text{i.e.}\quad\uline{u}(\rho)=(\varepsilon_2g).\rho
\tag{78}
\]

\uline{Corollaire} Soit $u$ un autom \unsure{du} groupe
à lacets $\widehat{\Pi}_{0,3}$, \unsure{commutant} extérieurement à
$\unsure{\mathfrak S_3|\mathcal S_3}$ (i.e. à $\sigma_\infty$ et à $\rho$), $\mu$
son multiplicateur, désignons aussi par
$u$ l'extension canonique de $u$ à
$\unsure{\widehat{\mathfrak S}_{0,3}^{+}|\widehat{\mathcal S}_{0,3}^{+}}
\simeq\mathrm{Sl}(2,\mathbb Z)^{\wedge\prime}$,
induisant l'identité sur
$\unsure{\mathfrak S_3|\mathcal S_3}
\simeq\unsure{\widehat{\mathfrak S}_{0,3}^{+}|\widehat{\mathcal S}_{0,3}^{+}}/
\widehat{\Pi}_{0,3}$,
et considérons le quotient $\mathrm{Sl}(2,\widehat{\mathbb Z})'$ de
$\mathrm{Sl}(2,\mathbb Z)^{\wedge\prime}$. Alors l'action de $u$ passe
en une action dans le quotient $\mathrm{Sl}(2,\widehat{\mathbb Z})'$,
et cette dernière est induite par autom.
\deleted{\unsure{intérieur}}\inserted{"\unsure{intérieur|intérieurs}"}
de $\mathrm{Gl}(2,\widehat{\mathbb Z})$, définie par une
\inserted{matrice} $\uline{u}\in\mathrm{Gl}(2,\widehat{\mathbb Z})$
telle que $\det u=\mu$
\margin{157}

\scanpage{159}
\mspage{369}

(ce qui décrit l'action \uline{extérieure} de $u$ sur
$\mathrm{Sl}(2,\widehat{\mathbb Z})'$ en termes du multiplicateur $\mu$).

\margin{\uline{NB} Pour vérifier le \unsure{corollaire}, on est
ramené au cas où $u$ commute à $\sigma_\infty$,
quitte à conjuguer avec un autom. intérieur \unsure{de}\illegible{1}\ldots}
\margin{On a fait ceci, pour le moment, pour $u$
commutant à $\sigma_\infty$.}

En fait, on a beaucoup mieux encore,
puisqu'on va définir \uline{canoniquement}
\inserted{\uline{modulo signe}},
en termes de $u$, la matrice $\uline{u}$ par la
formule (75), où \deleted{$t\text{ et }u\in$} $t,u\in\widehat{\mathbb Z}$ sont
définis ainsi : \deleted{\unsure{On a}} la donnée de l'autom $u$
définit $f\in\widehat{\Pi}_{0,3}$, modulo opération $f\mapsto\ell_0^n f$
(pour $n\in\widehat{\mathbb Z}$), et on désigne encore par $f$
l'image de $f$ dans $\mathrm{Sl}(2,\widehat{\mathbb Z})/\pm1$, qui est donc
une matrice $\begin{pmatrix}r&s\\t&u\end{pmatrix}$, définie mod signe,
On notera que dans $\mathrm{Sl}(2,\widehat{\mathbb Z})$ on a
\[
\ell_0^n\begin{pmatrix}r&s\\t&u\end{pmatrix}
=\begin{pmatrix}1&-2n\\0&1\end{pmatrix}
\begin{pmatrix}r&s\\t&u\end{pmatrix}
=\begin{pmatrix}r-2nt&s-2nu\\t&u\end{pmatrix}
\tag{79}
\]
($n\in\widehat{\mathbb Z}$), donc $(t,u)$ ne change pas (mod signe)
quand on remplace $f$ par $\ell_0^n f$ --- donc $(t,u)$
est déterminé mod signe par l'autom $u$
de $\widehat{\Pi}_{0,3}$. \deleted{\illegible{1}} L'observation analogue s'applique
d'ailleurs à $\alpha_0=\begin{pmatrix}a&b\\c&d\end{pmatrix}$, qui est déterminé
mod signe et mod multiplication à gauche
par un $\ell_0^n$ --- le \deleted{\illegible{1}} vecteur $(c,d)$ est
déterminé mod signe, donc \unsure{l'élément}
\[
c\rho+d
\]
\margin{158}

\scanpage{160}

est déterminé mod \uline{signe}, et \deleted{\illegible{1}}\inserted{sauf erreur}
on doit avoir une relation
\[
\tag{80}
\uline{u}\overset{?}{=}\mu\alpha(c\rho+d)
=\mu f^{-1}\ell_0^{-\nu}\alpha_0(c\rho+d)
\]
(cf (69) (73)), tout au moins si $\varepsilon_2=1$ ---
sinon il y a lieu peut-être de corriger par un facteur $\tau_\infty$ quelque part.
\deleted{Or je n'ai pas \unsure{envie} pour le moment de reprendre \unsure{mes|mon} \unsure{calculs|calcul} dans ce sens.}

\margin{NB cette formule est \unsure{valable} \unsure{moyennant} le facteur correctif $\delta=\varepsilon_3$, cf (82) plus bas\ldots}

\[
\ell_0^{-\nu}\alpha_0(c\rho+d)
=
\begin{pmatrix}1&2\nu=\mu-1\\0&1\end{pmatrix}
\underbrace{
\begin{pmatrix}a&b\\c&d\end{pmatrix}
\begin{pmatrix}d&c\deleted{\unsure{+d}}\\-c&c+d\end{pmatrix}
}_{\displaystyle
\underbrace{
\begin{pmatrix}
ad-bc=\delta&ac+bc+bd=\delta B\\
0&c^2+cd+d^2=\delta D
\end{pmatrix}
}_{\displaystyle\delta\begin{pmatrix}1&B\\0&D\end{pmatrix}}
}
\]
\[
\ell_0^{-\nu}\alpha_0(c\rho+d)
=\delta\begin{pmatrix}1&B+2\nu D=C\\0&D\end{pmatrix}
\]
i.e.
\[
\tag{81}
\ell_0^{-\nu}\alpha_0(c\rho+d)
=\delta\begin{pmatrix}1&C\\0&D\end{pmatrix}
\]
d'où
\[
\begin{aligned}
\mu\alpha(c\rho+d)
&=(\mu f^{-1})(\ell_0^{-\nu}\alpha_0(c\rho+d))
=\mu\delta\begin{pmatrix}u&-s\\-t&r\end{pmatrix}
\begin{pmatrix}1&C\\0&D\end{pmatrix}\\
&=\mu\delta\begin{pmatrix}u&uC-sD\\-t&-tC+rD\end{pmatrix}\\
&=\delta\begin{pmatrix}\mu u&-s+\mu uC\\-\mu t&r-\mu tC\end{pmatrix}
\end{aligned}
\]
Or (61)
\[
C=\varepsilon_2T=\varepsilon_2(rt+su)\qquad\text{d'où}
\]
\[
\begin{aligned}
-s+\mu uC
&=-s+\varepsilon_2\unsure{\mu urt|\mu turt}+\varepsilon_2\mu u^2s\\
&\overset{\inserted{\unsure{\text{si }\varepsilon_2=+1}}}{=}
 s\bigl(\underbrace{-1+\unsure{\varepsilon_2u^2|\varepsilon_2\mu u^2}}_{-\varepsilon_2\mu t^2}\bigr)
 +\varepsilon_2\mu urt\\
&=\varepsilon_2\mu t\bigl(\overbrace{\unsure{-st+ur|st+ur}}^{+1}\bigr)\\
&=\varepsilon_2\mu t
\end{aligned}
\]

\margin{159}

\scanpage{161}
\mspage{370}

\[
\begin{aligned}
r-\mu tC
&\overset{\inserted{\unsure{\text{si }\varepsilon_2=1}}}{=}
\underbrace{r-\varepsilon_2\mu rt^2}_{r(1-\varepsilon_2\mu t^2)=\varepsilon_2r\mu u^2}
-\varepsilon_2\mu stu\\
&=\varepsilon_2\mu u(ru-st)=\varepsilon_2\mu u,
\end{aligned}
\]
donc
\[
\mu\alpha(c\rho+d)
=\delta\mu\begin{pmatrix}u&\varepsilon_2\unsure{t}\\-t&\varepsilon_2u\end{pmatrix}
\overset{\deleted{\unsure{\text{si }\varepsilon_2=1}}}{=}\delta\uline{u},
\]
i.e.
\[
\tag{82}
\uline{u}=(\unsure{\delta}\mu)\alpha(c\rho+d)
=(\delta\mu)(f^{-1}\ell_0^{-\nu}\alpha_0)(c\rho+d),
\]
\deleted{\unsure{si $\varepsilon_2=1$ (i.e. $\mu\equiv1\ (4)$)}}
\inserted{où $\delta=ad-bc\in\{\pm1\}$, $\delta=\varepsilon_3$
(donc $\delta\equiv\mu\ (3)$)}.

Formule qui s'écrit aussi
\deleted{i.e.} \unsure{compte} \deleted{\illegible{1}} tenu de (75)
\[
\tag{83}
\left\{\begin{aligned}
\deleted{\unsure{f}}(t\sigma_\infty+u)\inserted{\tau_\infty^{\alpha_2}}
&=\delta\alpha(c\rho+d)
\qquad\deleted{\unsure{\text{si }\varepsilon_2=1}}\\
\text{où}\quad\alpha&=f^{-1}\ell_0^{-\nu}\alpha_0,
\quad\deleted{\illegible{1}}\quad
\delta=\varepsilon_3\in\{\pm1\},\quad\delta=ad-bc,\quad\delta\equiv\mu\ (3)
\end{aligned}\right.
\]
ou encore
\[
\tag{84}
f(t\sigma_\infty+u)\tau_\infty^{\alpha_2}
=\delta\ell_0^{-\nu}\alpha_0(c\rho+d)
\]
\[
\begin{array}{c|c}
\underbrace{
\begin{pmatrix}r&s\\t&u\end{pmatrix}
\begin{pmatrix}u&t\\-t&u\end{pmatrix}
}_{\displaystyle\begin{pmatrix}1&rt+su\\0&t^2+u^2\end{pmatrix}}
\begin{pmatrix}1&0\\0&\varepsilon_2\end{pmatrix}
&
\underbrace{
\begin{pmatrix}1&2\nu\\0&1\end{pmatrix}
\begin{pmatrix}a&b\\c&\unsure{d}\end{pmatrix}
\begin{pmatrix}d&c\\-c&c+d\end{pmatrix}
}_{\displaystyle\begin{pmatrix}1&C\\0&D\end{pmatrix}}
\quad\text{cf (81)}\\
\begin{pmatrix}1&rt+su\\0&t^2+u^2\end{pmatrix}\deleted{\begin{pmatrix}\unsure{u}&\unsure{t}\\-\unsure{t}&\unsure{u}\end{pmatrix}}&\\
\deleted{=\begin{pmatrix}1&T\\0&U\end{pmatrix}}&\\
\deleted{\unsure{=}}&\\
=\begin{pmatrix}1&\varepsilon_2T\\0&\varepsilon_2U\end{pmatrix}&
\end{array}
\]
et sous cette forme la formule \unsure{était} en fait connue,
comme le contenu essentiel de la \deleted{formule} relation des lacets sous forme matricielle (56) ou (57)
\[
:\quad\varepsilon_2U=D,\qquad\varepsilon_2T=C\qquad!
\]

\uline{Remarque} En fait, sous une fausse apparence d'arithmétique, on a fait \unsure{par|pr} la bande des gammes d'algèbre linéaire, valables sur des

\margin{160}

\scanpage{162}

anneaux quelconques et \unsure{formulable|formulables} sans doute sous des conditions beaucoup plus générales et plus ``\unsure{compréhensives}'' --- concernant des équations de la forme
\[
[\alpha_0,\rho]=\xi
\]
où $\rho\in\unsure{\mathrm{Gl}}(2,A)$ est donné, \deleted{\unsure{et}\illegible{1}}\inserted{ainsi que} $\xi$, et
$\alpha_0\in\mathrm{Gl}(2,A)$ est l'inconnue (assujettie éventuellement à des conditions du type $\det\alpha_0=\delta$, $\delta$ fixé), ainsi que des équations simultanées en $\alpha_0$ et $f$
\[
\tag{*}
[\alpha_0,\rho]=\varepsilon.\unsure{\mathcal{E}_0|\varepsilon_0}(\mu-1)[f,\sigma]
\qquad\text{cf (56)}
\]
où $\rho,\sigma\in\mathrm{Gl}(2,A)$ sont donnés, ainsi qu'un scalaire $\varepsilon\in A^*$, et
$\lambda\longmapsto\unsure{\mathcal{E}_0|\varepsilon_0}(\lambda)$ est un ``sous-groupe à un paramètre'', p.ex. on peut supposer
\[
\unsure{\mathcal{E}_0|\varepsilon_0}(\lambda)=\begin{pmatrix}1&-\lambda\\0&1\end{pmatrix}=\varepsilon_0^\lambda,
\]
\[
\varepsilon_0=\begin{pmatrix}1&-1\\0&1\end{pmatrix}
\]
--- en supposant au besoin une relation de la forme $\sigma\rho=\varepsilon_0$.

Modulo des ajustements agréables, toutes les formules du présent \S\ devraient pouvoir se développer dans ce contexte plus général
\inserted{(ainsi que des compléments substantiels)}
(--- mais je ne sais si je vais expliciter

\margin{161}

\scanpage{163}
\mspage{380}

ici ces \unsure{gammes} \unsure{délectables} d'algèbre linéaire,
voire d'interprétations schématiques du genre que le schéma des solutions de (*)
(pour $f,\alpha_0,\mu$ variables) est un schéma \inserted{relatif} lisse de dim relative 4, \unsure{qui etc.}

Je vais seulement expliciter, dans le contexte présent adélique, \unsure{trop} particulier manifestement, la

\uline{Proposition} Soient $\deleted{\unsure{a,b}}c,d\quad\deleted{\unsure{r,s}}t,u\in\widehat{\mathbb Z}$,
pour qu'ils \unsure{correspondent} à des données
$\alpha_0=\begin{pmatrix}a&b\\c&d\end{pmatrix}$,
$f=\begin{pmatrix}r&s\\t&u\end{pmatrix}$ satisfaisant
\[
\underbrace{ad-bc}_{\substack{\unsure{\text{déf}}\\=\varepsilon_3}}\in\{\pm1\},
\quad ru-st=1,
\quad\alpha_0\rho(\alpha_0)^{-1}=\varepsilon_2\ell_0^{\nu}(f\sigma_\infty(\unsure{f})^{-1})
\]
avec $\varepsilon_2\in\{\pm1\}$ et $\unsure{\nu}\in\widehat{\mathbb Z}$
\unsure{convenable|convenables} (cf (48) (49) (53)),
il \unsure{f|faut} et \deleted{\illegible{1}} il suffit qu'on ait
\[
\tag{85}
c^2+\unsure{cd}+d^2=\pm(t^2+u^2)\quad\in\widehat{\mathbb Z}^*
\]
Alors \deleted{\unsure{signes} $\varepsilon_2,\varepsilon_3$ et le scalaire $\unsure{\nu}\in\widehat{\mathbb Z}$ sont}
\inserted{$(\varepsilon_2,\varepsilon_3,2\nu+1)\in\{\pm1\}\times\{\pm1\}\times\widehat{\mathbb Z}^*$ est déterminé}
modulo signe (i.e. modulo passage à : $(-\varepsilon_2,-\varepsilon_3,-\nu)$)
\deleted{alors déterminés de façon unique ainsi : On}
\inserted{par $a,b,c,d$, par la \deleted{condition} relation}
\inserted{(où $\mu=2\nu+1$)}
\deleted{définit \illegible{1}}
\[
\deleted{\varepsilon_3\unsure{\equiv|=}\unsure{(c^2+cd+d^2)}\unsure{\in\widehat{\mathbb Z}^*}}
\]
\[
\tag{86}
\frac1\mu=\varepsilon_2(t^2+u^2)=\varepsilon_3(c^2+cd+u^2)
\]
\margin{\unsure{par}}
\[
\deleted{\unsure{\mu|\frac1\mu}=\unsure{\varepsilon_3}(c^2+cd+d^2)=}
\]
\deleted{\illegible{3}}
\inserted{(qui détermine $\varepsilon_3,\mu$ quand on a choisi $\varepsilon_2$)}
et on a
\[
\tag{87}
\left\{\begin{aligned}
\varepsilon_2&\equiv\mu\quad(4)\\
\varepsilon_3&\equiv\mu\quad(3)\\
\nu&=(\mu-1)/2
\end{aligned}\right.
\]

\margin{162}

\scanpage{164}

Les solutions s'obtiennent alors ainsi : on choisit
$\varepsilon_2,\varepsilon_3,\unsure{\mu|\mu_2}$ comme dessus (deux possibilités), \deleted{puis} et arbitrairement $r,s\in\widehat{\mathbb Z}$ satisfaisant
\[
ur-ts=1
\]
(il y en a toujours (car $r^2+s^2\in\widehat{\mathbb Z}^*$ --- donc
$\forall p\in\mathbb P$, $r_p$ et $s_p$ ne sont pas tous deux dans $p\mathbb Z_p$)
et si $r_0,s_0$ est une solution, les autres s'obtiennent sous forme paramétrique en $\lambda$
\[
\tag{87}
(r,s)=(r_0,s_0)\unsure{+|-}\lambda(\unsure{t|u},\unsure{u|t})
=(r_0\unsure{+|-}\lambda t,s_0\unsure{+|-}\lambda u).
\]
Ayant choisi $r,s$ donc $f=\begin{pmatrix}r&s\\t&u\end{pmatrix}\in\mathrm{Sl}(2,\widehat{\mathbb Z})$,

\deleted{
on en déduit
\[
f\sigma_\infty(f)^{-1}=f\,{}^tf
=\begin{pmatrix}R&S\\T&U\end{pmatrix}
\]
où $R=\unsure{r}^2+s^2$, $S=T=rt+su$, $U=t^2+u^2$, et on détermine successivement les quantités
\[
\tag{88}
\left\{\begin{aligned}
D&=1/\mu=\varepsilon_3(c^2+cd+d^2)
=\deleted{\varepsilon_2\underbracket{\unsure{U}}_{\inserted{t^2+}}=\varepsilon_2(\unsure{t^2}+\illegible{1})}
\inserted{\varepsilon_2(t^2+u^2)=\varepsilon_2U}\\
C&=\varepsilon_2T\ (=\varepsilon_2S)\\
B&=C+D-1=\varepsilon_2(S+U)-1\\
A&=\frac1D(1+BC)=\frac1{\varepsilon_2U}(1+S(S+U)-\varepsilon_2S)
\end{aligned}\right.
\]
\margin{$\unsure{t^2+}$}
d'où une matrice $\begin{pmatrix}A&B\\C&D\end{pmatrix}$ satisfaisant
$C+D-B=1$ et $D\in\widehat{\mathbb Z}^*$, et satisfaisant la relation (57)
\[
\begin{pmatrix}A+2\nu C&B+2\nu D\\C&D\end{pmatrix}
=\varepsilon_2\begin{pmatrix}R&S\\T&U\end{pmatrix}
\]
qui se réduit ici \unsure{à} la relation du coin sup. gauche
$A+2\nu C=\varepsilon_2R$ i.e.
}

\margin{163}

\scanpage{165}
\mspage{381}

il reste à choisir $a,b$ de façon à satisfaire aux deux relations
$\det\begin{pmatrix}a&b\\c&d\end{pmatrix}=\varepsilon_3$
et $C=\varepsilon_2 T$ (\unsure{cf.} (61) (62))
\inserted{i.e.}
\[
\tag{88}
\left\{\begin{aligned}
d\uline{a}-c\uline{b}&=\varepsilon_3\\
(\unsure{c}+d)\uline{a}+d\uline{b}
&=\varepsilon_2\varepsilon_3
\underbrace{(\uline{r}t+\uline{s}u)}_{S}
+\underbrace{(c^2+cd+d^2)}_{1/\mu\varepsilon_3}
\end{aligned}\right.
\]
qui est un système de deux équations linéaires en $\uline{a},\uline{b}$,
dont le déterminant est comme par hasard
$\unsure{c^2}+cd+d^2=1/\mu\varepsilon_3\in\widehat{\mathbb Z}^{*}$
--- ce qui assure une solution unique !

Il semble plus naturel peut être de partir du multiplicateur $\mu$, et de
chercher les systèmes
$\begin{pmatrix}a&b\\c&d\end{pmatrix}$
$\begin{pmatrix}r&s\\t&u\end{pmatrix}$ qui y correspondent,
mais en mettant l'accent sur $(c,d)\ (t,u)$ en première ligne.
À l'aide de $\mu$ on détermine $\varepsilon_2,\varepsilon_3$ par
\[
\mu\equiv\varepsilon_2\ (2),\qquad
\mu\equiv\varepsilon_3\ (3),
\]
\deleted{\unsure{puis}} ce qui permet de \deleted{\unsure{déf.}} représenter
$1/\mu$ sous la forme
\[
1/\mu=\varepsilon_2(t^2+u^2)=\varepsilon_3(c^2+cd+d^2)
\]
(avec $t,u,c,d\in\widehat{\mathbb Z}$, les systèmes $(t,u)$ comme
$(c,d)$ ``\unsure{dépendent} de $1$ paramètre'', \inserted{chacun}
en un sens très concret qu'on pourrait sans doute préciser si on y tenait,
\deleted{\unsure{un grand choix}}\inserted{avec des}
paramètres \deleted{\illegible{1}} $\sigma,\tau$ explicites\ldots{} Une
fois $c,d,\unsure{t,u}$ choisis, on choisit $r,\unsure{s|t}$ comme il a

\margin{164}

\scanpage{166}

été dit (avec un paramètre $\unsure{\lambda}$ \inserted{intervenant}
linéairement), ce qui détermine ensuite $r,s$ de façon unique, d'une façon
qui fait intervenir $\lambda$ \deleted{linéairement} de façon
linéaire-affine.

En fait, pour $c,d,t,u$ \inserted{et $\varepsilon_2,\varepsilon_3$}
fixés, \deleted{les équations}\inserted{de façon à}
satisfaire : $\varepsilon_2(t^2+u^2)=\varepsilon_3(c^2+cd+d^2)$,
les trois équa. qui relient $a,b,r,s$ sont \unsure{linéaires-affines}
\[
\tag{89}
\left\{\begin{aligned}
d\uline{a}-c\uline{b}&=\varepsilon_3\\
u\uline{r}-t\uline{s}&=1\\
(c+d)\uline{a}+d\uline{b}
-\varepsilon_2\varepsilon_3(u\uline{s}+t\uline{r})&=c^2+cd+d^2
\end{aligned}\right.
\]
\margin{on a souligné les ``coefficients inconnus''
$\uline{a},\uline{b},\uline{r},\uline{s}$ à déterminer}
et le \deleted{déterminant de ce} système est ``de rang $3$''
\deleted{\unsure{det}}
--- l'ens. de ses solutions est un ``espace affine de dim $1$''
sur $\widehat{\mathbb Z}$\ldots

\uline{Remarque} On n'a pas pour le moment tenu compte des conditions de
congruence
\[
\begin{pmatrix}a&b\\c&d\end{pmatrix}\equiv1\ (2)
\qquad
\begin{pmatrix}r&s\\t&u\end{pmatrix}\equiv1\ (\unsure{2})
\]
qui expriment que ces matrices proviennent mod signes par ``réduction''
d'éléments de
$\unsure{\widehat{\widetilde{\Pi}}_{0,3}}$, $\widehat{\Pi}_{0,3}$.
Pour les quantités $(c,d)$ \deleted{\unsure{fixées}} on trouve les
conséquences
\[
\left\{\begin{aligned}
d&\equiv\pm1\ (4)\quad\text{donc }d^2\equiv1\ (4)\\
c&\equiv0,2\ (4)\quad\text{donc }c^2\equiv0\ (4)
\end{aligned}\right.
\ \Longrightarrow\
\begin{aligned}
c^2+d^2+cd&\equiv1+cd\ (4)\\
&\equiv1+c\ (4)
\end{aligned}
\]
donc
\[
c^2+d^2+cd\equiv1+c\ \deleted{\unsure{(4)}}(4)
\quad\deleted{\unsure{\Longleftrightarrow}\ \unsure{cd}\ cd\equiv0\ (4)}
\quad\text{donc}
\]
donc
\[
\mu=\varepsilon_3(c^2+cd+d^2)\equiv
\deleted{\unsure{\Longleftrightarrow c\equiv0\ (4)}}
\varepsilon_3(1+c)\ (4)
\]
\deleted{(NB $d$ est inv. mod $4$ puisque $d\equiv1\ (2)$),}

\margin{165}

\scanpage{167}
\mspage{382}

On voit donc que la connaissance de $c$ \deleted{\illegible{1}}
\inserted{mod $4$}, et celle du signe $ad-bc=\varepsilon_3$,
permet de trouver le signe $\varepsilon_2$ \deleted{\unsure{par}} par
\[
\deleted{\unsure{\varepsilon_2}\equiv\varepsilon_3\unsure{c}\quad
\varepsilon_2\equiv\varepsilon_3(1+c)\ (4)}
\]
\[
\tag{90}
\varepsilon_2\equiv\varepsilon_3(1+c)\ (4)
\quad\text{i.e.}\quad
\varepsilon_2\varepsilon_3\equiv1+c\ (4)
\]
Si on se propose, pour $c,d,t,u$ fixés, satisfaisant
\[
\tag{91}
\left\{\begin{aligned}
c&\equiv t\equiv0\ (2)\qquad d\equiv u\equiv1\ (2)\\
(c^2+\unsure{cd}+d^2)&\unsure{\equiv|=}\deleted{\unsure{(1+c}}
\varepsilon(t^2+u^2)\in\widehat{\mathbb Z}^{*}
\end{aligned}\right.
\]
où $\varepsilon$ \deleted{\illegible{1}}\inserted{$\in\{\pm1\}$}
est le signe défini par
\[
\varepsilon\equiv1+c\ (4)
\]
($\varepsilon=1$ si $c\equiv0\ (4)$, $\varepsilon=-1$ si
$c\equiv2\ (4)$), de trouver les $a,b,r,s$ correspondants,
satisfaisant les conditions (89), et de plus les conditions de congruence
\[
\tag{92}
\unsure{b}\equiv\unsure{s}\equiv0\ (2)\qquad
a\equiv r\equiv1\ (2),
\]
on est amené à relever une solution déjà donnée mod $2$ de ce système
d'équations en une solution. En termes du \inserted{$\widehat{\mathbb Z}$-}
torseur $\unsure{\mathcal T}$ des solutions, cela signifie qu'on se donne
un élément du $\mathbb Z/2\mathbb Z$-torseur quotient correspondant
--- l'ens. de ses relèvements s'identifie à un torseur sous
$2\widehat{\mathbb Z}\simeq\widehat{\mathbb Z}$\ldots{} Aucun mystère !

\margin{166}

\scanpage{168}

\uline{Remarque} Le groupe additif $\widehat{\mathbb Z}$ opère sur l'ens.
des solutions
$(\alpha_0,f)=\left(\begin{pmatrix}a&b\\c&d\end{pmatrix},
\begin{pmatrix}r&s\\t&u\end{pmatrix}\right)$
des équations envisagées (48)(49)(56)
\[
\tag{93}
\det\alpha_0\in\{\pm1\},\quad\det f=1,\quad
\alpha_0\rho(\alpha_0)^{-1}
=\varepsilon\ell_0^{\nu}\bigl(f\sigma_\infty(f)^{-1}\bigr),
\quad\varepsilon\in\{\pm1\},
\]
par
\[
\tag{\unsure{93|94}}
\begin{aligned}
(\alpha_0,f)&\longmapsto(\alpha'_0,f')
=(\ell_0^n\alpha_0,\ell_0^nf)\quad\deleted{=(\unsure{a}-2nc,\ b-}\\
&=\left(\begin{pmatrix}a-2nc&b-2nd\\c&d\end{pmatrix},
\begin{pmatrix}r-2nt&s-2nu\\t&u\end{pmatrix}\right)
\end{aligned}
\]
($n\in\widehat{\mathbb Z}$).
Au facteur $2$ près dans $2n$, c'est justement l'indétermination dans
\deleted{\illegible{1}} l'ensemble des solutions \inserted{de (93)}, lorsque
$c,d,t,u$ sont fixés (cf (87) et proposition qui va avec).
\inserted{(le facteur $2$ est \unsure{assuré} \unsure{seulement quand}
\unsure{on considère des|en considérant des} solutions satisfaisant
$\alpha'_0\equiv f'_0\equiv\mathrm{Id}\ (2)$)}
On peut \deleted{\unsure{donc}} dire que ce qui est déterminé
intrinsèquement \inserted{dans
$\unsure{\mathrm{Gl}(2,\widehat{\mathbb Z})^{\pm1\prime}|\mathrm{Gl}(2,\widehat{\mathbb Z})^{\pm1}}$}
en termes de l'autom \inserted{à lacets} $\uline{u}$ de
$\widehat{\Pi}_{0,3}$ commutant à $\sigma_\infty$, c'est ni plus ni moins
que \deleted{les \unsure{valeurs}}
\[
\tag{95}
\begin{aligned}
\text{a)}\quad&\text{le multiplicateur }\mu\in\widehat{\mathbb Z}^{*}
\end{aligned}
\]
\deleted{\unsure{qui donne} $\varepsilon_2=\unsure{\varepsilon_2|\xi_2}(\mu),\
\varepsilon_3=\unsure{\varepsilon_3|\xi_3}(\mu)$}
\inserted{(\unsure{donc|d'où} $\unsure{\varepsilon_2}=\unsure{\varepsilon_2|\xi_2}(\mu),\
\varepsilon_3=\unsure{\varepsilon_3|\xi_3}(\mu)$)}

b) les quantités
\deleted{$\unsure{\delta}=\unsure{\varepsilon_3},\mu,\
c,d,t,u\in\widehat{\mathbb Z}$, \unsure{satisfaisant}}
les $c,d,t,u\in\widehat{\mathbb Z}$ satisfaisant la
sempiternelle relation (86)
\[
\tag{96}
\varepsilon_2(t^2+u^2)=\varepsilon_2(c^2+cd+d^2)=1/\mu
\]
et les relations de congruence
\[
\tag{97}
c\equiv t\equiv0\ (2)\qquad d\equiv u\equiv1\ (2)
\]

Il est vrai que l'on a vu qu'il y a un moyen \inserted{naturel} dans
$\widehat{\Pi}_{0,3}$ de normaliser le choix des $f$, de façon à lever
l'\unsure{ambiguité|ambiguïté} des $\ell_0^nf$
(cf \S\,41 \inserted{(24bis)} où c'était effleuré --- j'y reviendrai
encore par la suite\ldots) mais il me semble \inserted{\uline{peu}}
intéressant

\margin{167}
